% ID: SEL-D0003-OA-E02
% ESTADO: validada. S004, 2026-11-07; tres pasadas del mismo autor y revisión matemática, didáctica y visual.
\textbf{a) Dominio y reglas de derivación.}\par\nopagebreak[4]
Para $f$, el argumento del logaritmo debe ser positivo:
\[
\frac{1-x}{1+x}>0\iff -1<x<1.
\]
En este intervalo $1-x$ y $1+x$ son positivos y podemos escribir
\[
f(x)=\frac12[\ln(1-x)-\ln(1+x)].
\]
Aplicando la regla de la cadena,
\[
f'(x)=\frac12\left(-\frac1{1-x}-\frac1{1+x}\right)
=\boxed{\frac1{x^2-1}},\qquad -1<x<1.
\]
Aunque la expresión de la derivada tenga sentido en otros puntos, no amplía el dominio real de la función original.

Para $g$, definida en $\mathbb R$, usamos la regla del producto y la cadena:
\[
\begin{aligned}
g'(x)&=4x(x^2+1)e^{2x-1}+2(x^2+1)^2e^{2x-1}\\
&=2(x^2+1)e^{2x-1}[2x+x^2+1]\\
&=\boxed{2(x^2+1)(x+1)^2e^{2x-1}}.
\end{aligned}
\]

\textbf{b) Tangente y representación.}\par\nopagebreak[4]
La tangente en $x_0$ pasa por $(x_0,h(x_0))$ y tiene pendiente $h'(x_0)$. Aquí
\[
h(4)=16-24+8=0,\quad h'(x)=2x-6,\quad h'(4)=2.
\]
Por tanto, $y-h(4)=h'(4)(x-4)$ da \boxed{y=2x-8}.

La parábola $h(x)=(x-3)^2-1$ tiene vértice $(3,-1)$, eje $x=3$, cortes horizontales $(2,0)$ y $(4,0)$, y corte vertical $(0,8)$. Abre hacia arriba. La recta pasa por $(0,-8)$ y $(4,0)$ y toca la curva en este último punto.
\begin{center}\includegraphics[width=0.70\linewidth]{figuras/s004-tangente.pdf}\end{center}
Comprobación algebraica: $h(x)-(2x-8)=(x-4)^2\geq0$, con igualdad sólo en $x=4$; además, las dos pendientes allí son 2.
